%%%%%%%%%%%%%%%%%%%%%%%%%%%%%%%%%%%%%%%%%%%%%%%%%%%%%%%%%%%%%%%%%%%
%                                                                 %
%  GEANT manual in LaTeX form                                     %
%                                                                 %
%  Version 1.00                                                   %
%                                                                 %
%  Last Mod. 8 June 1993  17:20  MG                               %
%                                                                 %
%%%%%%%%%%%%%%%%%%%%%%%%%%%%%%%%%%%%%%%%%%%%%%%%%%%%%%%%%%%%%%%%%%%
\documentstyle[11pt,fleqn,epsfig,crngeant,bibunits,fontcmds]{cernman}
\newcommand{\Title}{GEANT User's Guide}%           Title for document
\psfigdriver{DVIPS}
\makeindex
\romanfont{times}
\PScommands% Initialize PS boxes
\newmathalphabet*{\mathtt}{cmtt}{m}{n}
\newmathalphabet*{\mathbf}{cmr}{b}{n}
\begin{document}
%  ==================== Front material ============================
%%%%%%%%%%%%%%%%%%%%%%%%%%%%%%%%%%%%%%%%%%%%%%%%%%%%%%%%%%%%%%%%%%%
%                                                                 %
%   GEANT -- Short write ups -- LaTeX Source                      %
%                                                                 %
%   Front Material: Title page,                                   %
%                   Copyright Notice                              %
%                   Preliminary Remarks                           %
%                   Table of Contents                             %
%   EPS file      : cern15.eps, cnastit.eps                       %
%                                                                 %
%   Editor: Michel Goossens / CN-AS                               %
%   Last Mod.: 11 June 1993 10:30 mg                              %
%                                                                 %
%%%%%%%%%%%%%%%%%%%%%%%%%%%%%%%%%%%%%%%%%%%%%%%%%%%%%%%%%%%%%%%%%%%
 
%%%%%%%%%%%%%%%%%%%%%%%%%%%%%%%%%%%%%%%%%%%%%%%%%%%%%%%%%%%%%%%%%%%%
%    Tile page                                                     %
%%%%%%%%%%%%%%%%%%%%%%%%%%%%%%%%%%%%%%%%%%%%%%%%%%%%%%%%%%%%%%%%%%%%
\latex{\def\Ptitle#1{\special{ps: /Printstring (#1) def}
                       \epsfbox{eps/cnastit.eps}}}
 
\begin{titlepage}
%begin{latexonly}
\vspace*{-23mm}%
\includegraphics[height=30mm]{cern.eps}
\hfill
\raise8mm\hbox{\Large\bf CERN Program Library Long Writeup W5013}
\hfill\mbox{}
\begin{center}
\mbox{}\\[6mm]
\mbox{\Ptitle{GEANT}}\\[3cm]
{\Huge Detector Description and}\\[2cm]
{\Huge Simulation Tool}\\[4cm]
{\Large Application Software Group}\\[6mm]
{\Large Computing and Networks Division}\\[2cm]
\end{center}
\vfill
\begin{center}\Large CERN Geneva, Switzerland\end{center}
%end{latexonly}
\begin{htmlonly}
\begin{center}
\Large CERN Program Library Long Writeup W5013\\[1cm]
\Huge GEANT\\[2cm]
\Large Detector Description and Simulation Tool\\[1cm]
\large IT/ASD Group\\
\large CERN, Geneva, Switzerland
\end{center}
\end{htmlonly}
\end{titlepage}

%%%%%%%%%%%%%%%%%%%%%%%%%%%%%%%%%%%%%%%%%%%%%%%%%%%%%%%%%%%%%%%%%%%%
%    Copyright  page                                               %
%%%%%%%%%%%%%%%%%%%%%%%%%%%%%%%%%%%%%%%%%%%%%%%%%%%%%%%%%%%%%%%%%%%%

\thispagestyle{empty}
\framebox[\textwidth][t]{\hfill\begin{minipage}{0.96\textwidth}%
\vspace*{3mm}
\begin{center}Copyright Notice\end{center}
\parskip\baselineskip

\textbf{GEANT -- Detector Description and Simulation Tool}
 
\copyright{} Copyright CERN, Geneva 1993
 
Copyright and any other appropriate legal protection of these
computer programs and associated documentation reserved in all
countries of the world.
 
These programs or documentation may not be reproduced by any
method without prior written consent of the Director-General
of CERN or his delegate.
 
Permission for the usage of any programs described herein is
granted apriori to those scientific institutes associated with
the CERN experimental program or with whom CERN has concluded
a scientific collaboration agreement.
 
CERN welcomes comments concerning the Geant code
but undertakes no obligation for the maintenance of the programs,
nor responsibility for their correctness, and accepts no liability
whatsoever resulting from the use of its programs.
 
Requests for information should be addressed to:
\vspace*{-.5\baselineskip}
\begin{center}\ttfamily
\begin{tabular}{l}
CERN Program Library Office              \\
CERN-IT Division                         \\
CH-1211 Geneva 23                        \\
Switzerland                              \\
Tel.      +41 22 767 4951                \\
Fax.      +41 22 767 7155                \\
Email:    cernlib@cern.ch                \\
WWW:      http://wwwinfo.cern.ch/asd/index.html
\end{tabular}
\end{center}
\vspace*{2mm}
\end{minipage}\hfill}%end of minipage in framebox
\vspace{6mm}
 
{\bf Trademark notice: All trademarks appearing in this guide are acknowledged as such.}
\vfill

\begin{tabular}{l@{\quad}l@{\quad}>{\small\ttfamily}l}
\emph{Contact Persons}:       & IT/ASD/SImulation section & (giani\atsign cern.ch)\\[1mm]
\emph{Documentation Consultant}: & Michel Goossens /IT    & (goossens\atsign cern.ch)\\[1cm]
\textem{Edition -- March 1994}
\end{tabular}
\newpage
 

%  ==================== Body of text ==============================
\pagenumbering{arabic}
\setcounter{page}{1}
 
%%%%%%   Catalog of Program packages and entries%%%%
 
\def\Rtnr{Catalog}%Dummy routine name to appear at bottom of page
%%%%%%%%%%%%%%%%%%%%%%%%%%%%%%%%%%%%%%%%%%%%%%%%%%%%%%%%%%%%%%%%%%%
%                                                                 %
%  GEANT manual in LaTeX form   	                          %
%                                                                 %
%  Michel Goossens (for translation into LaTeX)                   %
%  Version 2.00                                                   %
%  Last Mod.  2 Jun 1998  17:10   MG                              %
%                                                                 %
%%%%%%%%%%%%%%%%%%%%%%%%%%%%%%%%%%%%%%%%%%%%%%%%%%%%%%%%%%%%%%%%%%%
%%% Generated from headings
\chapter{Catalog of Geant sections}
\begin{DLtt}{12345678}
\item[AAAA001] Foreword
\item[AAAA002] Introduction to the manual
\item[BASE001] Introduction to GEANT
\item[BASE010] Simplified Program Flow Chart
\item[BASE020] The data structures and their relationship
\item[BASE040] Summary of Data Records
\item[BASE090] The reference systems and physical units
\item[BASE100] Examples of GEANT application
\item[BASE110] The system initialisation routines
\item[BASE200] Steering routines for event processing
\item[BASE280] Storing and retrieving JRUNG and JHEAD information
\item[BASE299] The banks JRUNG and JHEAD
\item[BASE300] Example of user termination routine
\item[BASE400] Debugging facilities
\item[BASE410] Utility Routines
\item[BASE420] The random number generator
\item[CONS001] Introduction to the section CONS
\item[CONS100] Material definition
\item[CONS101] Retrieve material cross-sections and stopping power
\item[CONS110] Mixtures and Compounds
\item[CONS199] The Material data structure JMATE
\item[CONS200] Tracking medium parameters
\item[CONS210] Special Tracking Parameters
\item[CONS300] Particle definition
\item[CONS310] Branching Ratios and Particle Decay Modes
\item[DRAW001] Introduction to the Drawing package
\item[DRAW010] The Ray-tracing package
\item[DRAW110] Drawing a Volume -- Case 1
\item[DRAW115] Drawing a Volume Projection view -- Case 2
\item[DRAW120] Draw a volume cut view
\item[DRAW130] Draw Particle Trajectories
\item[DRAW140] Drawing Track Hits in Sensitive Detectors
\item[DRAW210] Drawing the geometrical tree
\item[DRAW220] Drawing volume specifications
\item[DRAW300] Handling View banks
\item[DRAW399] The data structure JDRAW
\item[DRAW400] Utility routines of the drawing package
\item[GEOM001] The geometry package
\item[GEOM010] Tracking inside volumes and optimisation
\item[GEOM020] ``MANY'' Volumes and boolean operations on volumes
\item[GEOM050] The GEANT shapes
\item[GEOM100] Creation of a volume
\item[GEOM110] Positioning a volume inside its mother
\item[GEOM120] Positioning a volume inside its mother with parameters
\item[GEOM130] Division of a volume into a given number of cells
\item[GEOM140] Division of a Volume into cells of a given size
\item[GEOM150] Division of a volume - general case
\item[GEOM199] The volume data structure -- JVOLUM
\item[GEOM200] Rotation matrices
\item[GEOM299] The rotation matrix data structure JROTM
\item[GEOM300] Finding in which volume a point is
\item[GEOM310] Finding distance to next boundary
\item[GEOM320] Reference system transformations
\item[GEOM400] Pseudo-division of a volume
\item[GEOM410] Ordering the contents of a volume
\item[GEOM500] Volume attributes
\item[GEOM600] User initialisation of the common block /GCVOLU/
\item[GEOM700] Medium search statistics
\item[GEOM900] End of geometry definition
\item[GEOM910] The CADINT Interface
\item[HITS001] The detector response package
\item[HITS100] Sensitive DETector definition
\item[HITS105] Detector aliases
\item[HITS110] DETector hit parameters
\item[HITS120] DETector Digitisation parameters
\item[HITS130] User detector parameters
\item[HITS199] The SET data structure JSET
\item[HITS200] Routines to store and retrieve HITS
\item[HITS299] The JHITS data structure
\item[HITS300] Routines to store and retrieve DIGItisations
\item[HITS399] The JDIGI data structure
\item[HITS400] Intersection of a track with a cylinder or a plane
\item[HITS500] Digitisation for drift- or MWP- Chambers
\item[HITS510] Digitisation for drift chambers
\item[IOPA001] The I/O routines
\item[IOPA200] ZEBRA sequential files handling
\item[IOPA300] Data structure I/O with sequential files
\item[IOPA400] ZEBRA direct access files handling
\item[IOPA500] Data structure I/O with direct access files
\item[KINE001] Section KINE
\item[KINE100] Storing and retrieving vertex and track parameters
\item[KINE199] The data structures JVERTX and JKINE
\item[KINE200] Interface to the Lund Monte Carlo
\item[KINE210] $\tau ^{\pm }$ generation and decay
\item[PHYS001] Introduction to the section PHYS
\item[PHYS010] Compute the occurrence of a process
\item[PHYS100] Steering routine for physics initialisation
\item[PHYS210] Total cross-section for e+e- pair production by photons
\item[PHYS211] Simulation of e-e+ pair production by photons
\item[PHYS220] Total cross-section for Compton scattering
\item[PHYS221] Simulation of Compton scattering
\item[PHYS230] Total cross-section for photoelectric effect
\item[PHYS231] Simulation of photoelectric Effect
\item[PHYS240] Photon-induced fission on heavy materials
\item[PHYS250] Total cross-section for Rayleigh scattering
\item[PHYS251] Simulation of Rayleigh scattering
\item[PHYS260] \v{C}erenkov photons
\item[PHYS320] Gaussian multiple scattering
\item[PHYS325] Moli\`ere scattering
\item[PHYS328] Plural scattering
\item[PHYS330] Ionisation processes induced by e+/e-
\item[PHYS331] Simulation of the delta-ray production
\item[PHYS332] Simulation of energy loss straggling
\item[PHYS333] Information about energy loss fluctuations
\item[PHYS334] Models for energy loss fluctuations in thin layers
\item[PHYS337] Birks' saturation law
\item[PHYS340] Total cross-section and energy loss for bremsstrahlung by e-e+
\item[PHYS341] Simulation of discrete bremsstrahlung by electrons
\item[PHYS350] Total cross-section for e+e- annihilation
\item[PHYS351] Simulation of e+e- annihilation
\item[PHYS360] Synchrotron radiation
\item[PHYS400] Simulation of particle decays in flight
\item[PHYS410] Rotations and Lorentz transformation
\item[PHYS430] Ionisation processes for muons and protons
\item[PHYS431] Ionisation processes for heavy ions
\item[PHYS440] Total cross-section and energy loss for bremsstrahlung by Muons
\item[PHYS441] Simulation of discrete bremsstrahlung by muons
\item[PHYS450] Total cross-section and energy loss for e-e+ pair production by muons
\item[PHYS451] Simulation of e+e- pair production by muons
\item[PHYS460] Muon-nucleus interactions
\item[PHYS510] The GEANT/GHEISHA Interface
\item[PHYS520] The GEANT/FLUKA Interface
\item[PHYS530] The GEANT/MICAP interface
\item[TRAK001] The tracking package
\item[TRAK110] Steering routine to track one event
\item[TRAK120] Steering routine to track one particle
\item[TRAK130] Tracking one particle through a volume
\item[TRAK200] The tracking routines block diagrams
\item[TRAK300] Storing secondary tracks in the stack
\item[TRAK310] Altering the order of tracking secondary particles
\item[TRAK399] The temporary stack data structure JSTAK
\item[TRAK400] Handling of track space points
\item[TRAK499] The space point data structure JXYZ
\item[TRAK500] Tracking routines in magnetic field
\item[XINT001] The interactive version of GEANT
\item[XINT002] Introduction to the Interactive version of GEANT
\item[XINT010] Screen views of GEANT++
\item[ZZZZ010] List of COMMON Blocks
\item[ZZZZ999] Index of Documented GEANT routines
\end{DLtt}

 
\let\LARGE\large
\let\Large\large
\let\DL\DLtt 

% Here come the different files to be included
 
%%     TEMP part     %%
 
\begin{bibunit}[unsrt]
\renewcommand{\bibname}{TEMP Bibliography}
%%%%%%%%%%%%%%%%%%%%%%%%%%%%%%%%%%%%%%%%%%%%%%%%%%%%%%%%%%%%%%%%%%%
%                                                                 %
%  GEANT manual in LaTeX form                                     %
%                                                                 %
%  Version 1.00                                                   %
%  Last Mod.  8 June 1993  1300   MG                              %
%                                                                 %
%%%%%%%%%%%%%%%%%%%%%%%%%%%%%%%%%%%%%%%%%%%%%%%%%%%%%%%%%%%%%%%%%%%
\Authors {G.N. Patrick, L.Urban}   \Origin {Same}
\Version{Geant 3.16}\Routid{PHYS341}
\Submitted {26.10.84}  \Revised{19.12.92}
\Makehead{Simulation  of   discrete Bremsstrahlung by electrons}

\section{Subroutines}

\Shubr{GBREME}{}

\Rind{GBREME} generates a photon from Bremsstrahlung of an electron by
treating it as a discrete process.
The secondary photon energy is sampled from a parameterization of
the Bremsstrahlung calculation of Seltzer and Berger
\cite{bib-SEL1} for electron energies below 10 GeV, and from the screened
Bethe-Heitler cross-section above 10 GeV.
Midgal corrections are applied in both cases. 
For the angular distribution of the photon, is calls the function
\Rind{GBTETH}.
 
Input :  via common block \FCind{/GCTRAK/}
 
Output:  via common block \FCind{/GCKING/}.
 
\Rind{GBREME} is called automatically from the tracking
routine \Rind{GTELEC} if, and when,
the parent electron reaches its radiation point during the
tracking stage of GEANT.
 
\Sfunc{GBTETH}{THETA = GBTETH(ENER,PARTM,EFRAC)}

\Rind{GBTETH} calculates the angular distribution for the $e^+ e^-$-pair
in pair production and for the photon in Bremsstrahlung.
In case of Bremsstrahlung it gives the
scaled angle for a photon radiated by an electron (mass {\tt PARTM}) of 
energy {\tt ENER}. The energy of photon is {\tt EFRAC} times
the initial energy of the electron.
\Rind{GBTETH} is called by \Rind{GBREME}.

\section {Method}
 
The Coulomb corrected Bethe-Heitler differential cross-section for
production of a photon of energy $\epsilon E$ 
($\epsilon = k/E$, where $k$ is the photon energy)
by an electron of primary energy $E$ is given as in 
\notHTML{\cite{bib-WILL,bib-BUTC,bib-EGS3}}\HTML{\cite{bib-WILL}, \cite{bib-BUTC}, \cite{bib-EGS3}}:
 
\begin{equation}
\frac{d\sigma(Z,E,\epsilon)} {d \epsilon} =
      \frac{r_0^2 \alpha Z[Z +\xi(Z)]}{\epsilon}
    \left\{[1+(1 -\epsilon)^2][\Phi_1(\delta )
   - F(Z)] -\frac{2} {3}(1 - \epsilon)[\Phi_2(\delta) - F(Z)]\right\}
\end{equation}
where, $\Phi_i(\delta)$ are the screening functions depending on the screening
variable $\delta$
\[ \begin{array}{ll}

\begin{array}{l}
\delta  =  \frac {\Dstm 136m}{\Dstm Z^{1/3}E} \end{array}
            \frac{\Dstm \epsilon}{\Dstm (1-\epsilon)}
              &        m = \mbox{electron mass} \\
& \\
\noindent \left.\begin{array}{l}
\Phi_1(\delta)  =  20.867 - 3.242\delta +0.625\delta^2  \\
\Phi_2(\delta)  =  20.209 - 1.930\delta - 0.086\delta^2  \\
\end{array}        \right\} & \; \delta\leq 1 \\
\begin{array}{l} \Phi_1(\delta) = \Phi_2(\delta) =  
21.12 - 4.184 \ln(\delta+0.952) \end{array} 
                    & \; \delta > 1  \\
&\\
\begin{array}{l} F(Z) \end{array}  = \left\{\begin{array}{l}
        4/3 \ln Z          \\
        4/3 \ln Z+4f_c(Z)  \\
        \end{array}\right.      
&  \begin{array}{l}
E<0.05 \; GeV      \\
E\geq 0.05 \; GeV  \\
\end{array}  \\
\\
\begin{array}{l}  \xi(Z)=\frac{\Dstm\ln {(1440/Z^{2/3})}}{\Dstm\ln
{(183/Z^{1/3})}-f_c(Z)} \end{array}  & \\
\\
\begin{array}{l} f_c(Z) \mbox{\qquad the Coulomb correction function}
\end{array}  \\
\end{array} \]

%all these arrays of one column are here just to line up the text
\begin{eqnarray*}
f_c(Z) & = & a(1/(1+a)+0.20206-0.0369a+0.0083*a^2-0.002a^3)\\
    a   & = & (\alpha*Z)^2                                     \\
\alpha  & = & 1./137.
\end{eqnarray*}
 The kinematically allowed region for $\epsilon$ is
\begin{equation}
\epsilon_c \equiv \frac{k_c}{E} \leq \epsilon \leq  1- \frac{m}{E} 
\equiv \epsilon_1
\end{equation}
where,  $k_c$ is the photon cut-off energy below which the Bremsstrahlung
is treated as a continuous energy loss ({\tt BCUTE} in the program).
The cross-section (1) can be decomposed as
(the normalization is not important !)
\begin{equation}
\frac{d \sigma}{d \epsilon} =  \sum^{2}_{i=1} \alpha_i f_i(\epsilon)
      g_i(\epsilon )
\end{equation}
\begin{flushleft}
$\begin{array}{l@{}l@{\hspace{4cm}}l}
\mbox{where} & & \\
& \alpha_1 = F_{10}\ln \frac{\Tstm \epsilon_1}{\Tstm \epsilon_c}&
  \alpha_2 = F_{20}\frac{\Tstm 3}{\Tstm 8}\
             \frac{\Tstm \epsilon_1^2 -\epsilon_c^2}{\Tstm 1-\epsilon_c}\\
& & \\
& f_1 (\epsilon) = \frac{\Tstm 1}
                 {\Tstm \epsilon \; \ln {(\epsilon_1/ \epsilon_c)}} &
  f_2(\epsilon) = \frac{\Tstm 2 \epsilon}{\Tstm \epsilon_1^2 -\epsilon_c^2}\\
& g_1(\epsilon)=\frac{\Tstm F_1}{\Tstm F_{10}}
                \frac{\Tstm 1-\epsilon}{\Tstm 1-\epsilon_c} M(\epsilon)&
g_2(\epsilon )  = \frac{\Tstm F_2}{\Tstm F_{20}} . C_M (\epsilon )\\
\mbox{and} & & \\
& F_1(\delta) = 3\Phi_1(\delta)-\Phi_2(\delta)-2F(Z)  &
  F_{10}=F_1(\delta_{min})                     \\
& F_2(\delta) = 2\Phi_1(\delta)-2F(Z)                 &
  F_{20} = F_2(\delta_{min})                   \\
\\
& \delta_{min} = \frac{\Tstm 136m}{\Tstm Z^{1/3}E}
                \frac{\Tstm \epsilon_c}{\Tstm 1- \epsilon_c}
\end{array}$
\end{flushleft}
$\delta_{min}$ is the minimum value of the variable $\delta$, and
$C_M(\epsilon)$  is the Midgal correction factor \cite{bib-MIGD}:
\[
C_M (\epsilon)  =\frac{1 + C_0 / \epsilon_1^2}
               {1 + C_0 / \epsilon^2} 
\]
 where
\[
C_0 =\frac{nr_0 \lambdar^2 }{\pi}
\]
\begin{DL}{MM}
\item[$n$]           electron density in the medium
\item[$r_0$]         classical electron radius
\item[$\lambdar$]    reduced Compton wavelength of the electron.
\end{DL}
This correction decreases the cross-section for low photon energy.
Using the decomposition (3) the sampling of $\epsilon$ can be done by
(for a set of random numbers: $0\leq r_i\leq 1$)
\begin{OL}
\item  selecting the integer i (1 or 2)\\
$ \begin{array}{ll}
\mbox{if  }r_0 \leq BPAR = \alpha_1/(\alpha_1+\alpha_2), & \mbox{then  } i=1 \\
\mbox{if  } r_0 > BPAR,                                  & \mbox{then  } i=2
\end{array}$
\item
sampling $\epsilon$ from $f_i(\epsilon) $ \\
$\begin{array}{ll}
\epsilon =\epsilon c(\epsilon_1/\epsilon_c)^{r1} & \mbox{when  } i=1 \\
\epsilon  =\sqrt{\epsilon_c^2 + r_1 (\epsilon m^2-\epsilon_c^2)} 
                                                 & \mbox{when  } i=2 
\end{array}$
\item
calculating the rejection function $g_i(\epsilon)$ and\\
if $ r_2>g_i(\epsilon)$  rejecting $\epsilon$, starting again from 1\\
if $ r_2\leq g_i(\epsilon)$ accepting $\epsilon$
\end{OL}
As in the case of the pair production after sampling $\epsilon$ from
$f_i(\epsilon)$,
we have to check that
\[
\delta  = \frac{136 m} {Z^{1/3}E} \frac{\epsilon} {1 - \epsilon_c}
       \leq \delta_{max}=
       \exp  \left(\frac{21.12-f(Z)} {4.184}\right)  - 0.952
\]
and, if this is not the case, we have to start again from step 1.
 
The decomposition (3) used here is simpler and more effective
than that used in the earlier versions of GEANT and in {\tt EGS},
and the method outlined above has no convergence problems.
 
After the successful sampling of $\epsilon$, \Rind{GBREME} generates the
polar angles
of the radiated photon with respect to an axis defined along the parent
electron's momentum.
The energy-angle distribution is given by 
Tsai~\notHTML{\cite{bib-TSAI,bib-TSAI-err}}\HTML{\cite{bib-TSAI}, \cite{bib-TSAI-err}}
as following:
\begin{eqnarray}
\frac{d \sigma}{dkd \Omega} 
& = & \frac{2 \alpha^{2}e^{2}}{\pi k m^{4}}
      \left\{ \left[ \frac{2y-2}{(1+l)^2}+\frac{12l(1-y)}{(1+l)^4}\right]
      (Z^{2}+Z) \nonumber  \right. \\
&   & \mbox{} + \left. \left[ \frac{2-2y-y^{2}}{(1+l)^2}- 
      \frac{4l(1-y)}{(1+l)^4}
      \right]
      \left[ X-2Z^{2}f((\alpha Z)^{2})\right]
      \right\} \nonumber \\
\end{eqnarray}
where $k$ is the photon energy, $E$ is the initial electron energy,
$y = k/E$ and $l = E^{2} \theta^{2}/m^{2}$. This distribution
is quite complicated for sampling and, furthermore, for a variable
$u = E \theta/m$, shows a very weak dependence on $Z$, $E(k)$,
$y = k/E$. Thus, the distribution can be approximated by a function
\begin{equation}
f(u) = C \left( u e^{-au} + d u e^{-3au} \right)
\end{equation}
where
\begin{eqnarray*}
C & = & \frac{9a^{2}}{9 + d} \\
a & = & 0.625 \\
d & = & 27.0 \\
\end{eqnarray*}
The sampling of the function $f(u)$ can be done in the following way
($r_{i}$ are uniformly distributed random numbers 
in [0,1]):
\begin{enumerate}
\item Choose between $u e^{-au}$ and $d u e^{-3au}$, \\
      weights: $9/(9+d)$ and $d/(9+d)$. \\
      If $r_{1} < 9/(9+d)$ then $a=b$, \\
      else $b=3a$.
\item Sample $u e^{-bu}$, \\
      $u=-(\ln r_{2} + \ln r_{3})/b$.
\end{enumerate}
The azimuthal angle, $\Phi$, is generated isotropically.
This information is used to calculate the momentum vector of the radiated
photon, to transform it to the GEANT coordinate system and 
to store the result into common block \FCind{/GCKING/}. Also, the momentum of the parent electron 
is updated.

\subsection{Restrictions}

\begin{enumerate}
\item
Effects due to breakdown of the Born approximation at low energies are ignored.
\item
Target materials composed of compounds or mixtures are treated identically
to elements using the effective atomic number $Z$ calculated in
\Rind{GSMIXT}
(this is not the case when computing the mean free path!).
\end{enumerate}
\putbib[cnasbibl,geabibl]
\end{bibunit}

%  ==================== Index material ============================

\setcounter{page}{1}%                                Reset page counter
\def\Rtnr{Index}%Dummy routine name to appear at bottom of page
\input{\jobname.ind} % index

\end{document}
